\section*{数学中的真理概念。}
\phantomsection
\addcontentsline{toc}{section}{数学中的真理概念。}

数学家们一直确信，他们证明的是“真理”或“真命题”；这样的信念显然只能是感情上的或形而上学的，而且，站到数学的地盘上既不能为它辩护，也不能赋予它一个不沦为同义反复的意义。因此，数学中真理概念的历史属于哲学史而不属于数学史；但这一概念的演变对数学的演变有过无可否认的影响，因此我们不能对它保持沉默。

首先应当指出，要看到一位具有深厚哲学素养的数学家，与要看到一位通晓数学的哲学家同样罕见；数学家关于哲学问题的意见，即使这些问题涉及他们自己的领域，也多半是从可疑的来源二手、三手接受来的意见。但也恰恰因为如此，使数学史家感兴趣的正是这些平庸的意见，其程度至少不亚于对下述思想家独创见解的兴趣：如笛卡尔或莱布尼茨（这两人都是第一流的数学家）、柏拉图（他至少还跟得上他那个时代的数学）、亚里士多德或康德（对这后两位就不能这样说了）。

数学真理的传统概念可以一直追溯到文艺复兴时期。在这种观念里，数学家研究的对象与自然科学研究的对象之间没有多大差别；二者都是可知的，人都已经通过直觉和理性把握了它们；没有理由怀疑直觉或理性，它们只有在使用得当时才会出错。“必须”，帕斯卡说，“有一种全然错乱的精神，才会在如此重大、几乎不可能失手的原则上推理错误”（{[}244{]}，第 XII 卷，第~9~页）。笛卡尔在自家炉火旁确信，“历来只有数学家能够找到一些证明，也就是说一些确实可靠的理由”（{[}85 a{]}，第 VI 卷，第~19~页），而且（如果相信他的自述的话）这发生在他建成如下形而上学之前很久；在那形而上学中，“我先前当作规则的那件事，也就是我们能够设想得非常清楚、非常分明的一切对象都是真的这件事”，他说，“之所以得到保证，只是因为上帝存在，且他是一个完满的存在”（{[}85 a{]}，第 VI 卷，第~38~页）。如果说莱布尼茨反驳笛卡尔，说看不出如何才能认定一个观念是“清楚而分明的”，\(^{23}\) 那么他也认为，只要表述被理解了，公理就是定义的明显而无可争议的推论。\(^{24}\) 同样不应忘记，在当时的语言里，数学包含许多我们今天已不再承认其为科学的学问，有时甚至一直延伸到工程师的技艺；而在它们所激发的信心之中，它们应用于“自然哲学”、“机械技艺”和航海所取得的惊人成功占了很大一部分。

在这种看待事物的方式下，公理与演绎规则一样，不再容许被讨论或质疑；至多可以让每个人按照自己的偏好，去选择“以古人的方式”推理，还是任凭自己的直觉发挥。出发点的选择同样是个趣味问题，于是出现了为数众多的《Elements》“改写本”，其中《Elements》坚实的逻辑框架变成了一幅离奇的讽刺画；人们给出无穷小演算和理性力学的一些概述，号称是从一些建立得惊人糟糕的基本原理推演出来的；而斯宾诺莎或许也是出于诚意，把他的《伦理学》说成是“more geometrico demonstrata”（以几何方式证明的）。如果说在 17 世纪很难找到两位在任何问题上都意见一致的数学家，如果说论战日日发生、无休无止而尖刻激烈，真理的概念却始终不曾受到质疑。“关于每个对象只有一个真理”，笛卡尔说，“谁找到了它，谁就知道关于它所能知道的一切”（{[}85 a{]}，第 VI 卷，第~21~页）。

虽然希腊鼎盛时期关于这些问题的数学文本没有一份留存下来，但希腊数学家在这个问题上的观点很可能要细致得多。演绎规则只是凭借经验才被提炼到足以令人完全信服的地步；在它们可以被认为凌驾于一切讨论之前，人们必须经历的主要是摸索和悖理。而且，如果以为就连帕斯卡判定为最明显的那些“公理”（据他妹妹所传的说法，他童年时凭着万无一失的本能自己发现了它们）都不曾成为长期讨论的对象，那就误解了希腊人的批判精神，误解了他们对论辩和巧辩的嗜好。在一个严格说来并非几何学的领域里，埃利亚学派的悖论为我们保存了这类论战的一些痕迹；阿基米德在指出（{[}5 b{]}，第 II 卷，第~265~页）他的前辈们在许多情形使用了我们通常冠以他名字的那条公理时，补充说，借助这条公理所证明的东西“与不用它所证明的东西同样地被人们接受”，而他的结果只要在同样的基础上被接受就够了。柏拉图按照他的形而上学观点，把数学呈现为通达“自在真理”的途径，把数学处理的对象说成在理念世界中具有实在的存在；但他还是在《理想国》的一段著名文字中精确地刻画了数学方法：“凡是研究几何与算术的人 \ldots{} 都假定偶数与奇数、三种角；他们把它们当作已知的东西：一旦作了假定，他们就认为自己不再需要对自己或对别人说明它们，{[}认为它{]}对每个人都清楚；从那里出发，他们循序渐进，以便经由一致的同意达到他们的研究所提示的目标”（{[}250{]}，第 VI 卷，510 c-e）。因此，构成一个证明的，首先是一个提供任意起点的出发点（尽管“对每个人都清楚”），他说，在此之外人们不再试图深挖；然后是一场沿着一串中间阶段循序进行的讨论；最后，在每一步，都由对话者的同意来为推理的正确性作保。还必须补充的是：公理一经陈述，原则上便不再容许任何对直觉的新诉求；普罗克洛斯引述格米努斯时提醒说，“我们从这门科学的先驱者们那里学到，当涉及必须成为我们几何学说之一部分的推理时，不要把仅只貌似可信的结论放在眼里”（{[}153 e{]}，第 I 卷，第~203~页）。

因此，正是经验与批判的熔炉，造就了数学推理规则的提炼；而如果确实像有人以相当可信的方式论证过的那样 \([317\ d]\)，欧几里得《Elements》第八卷为我们保存了阿尔希塔斯算术的一部分，那么在那里看到那种相当迂腐的推理的生硬便不足为怪了——在任何一个“发现”了或自以为发现了“严格性”的数学学派里，这种生硬都不会不出现。但是，这些推理规则一旦进入数学家的实践，直到不久以前似乎从未受到过怀疑：在亚里士多德和斯多葛学派那里，虽然部分规则是通过推理格式从其他规则推出来的，初始规则却总被认为是自明的。同样，在回溯到那些在他们看来为他们时代的科学提供了坚实基础（例如他们在传统归于希俄斯的希波克拉底（约公元前 450 年）的第一部《Elements》中所提出的那些）的“假设”、“公理”和“公设”之后，古典时期的希腊数学家似乎把全部精力都投入到新结果的发现，而不是对这些基础的批判——在那个时候，这种批判不可能不是徒劳的；而且，撇开一切形而上学的考虑不谈，上文所引柏拉图的文本所见证的，正是数学家们在其科学基础上的这种普遍一致。

另一方面，希腊数学家似乎并不认为有可能解释清楚作为他们出发点的那些“初始概念”，如直线、曲面、量的比；他们如果给出其“定义”，那显然出于良心的不安，而对其效力不抱任何幻想。不言而喻，反过来，对于“初始概念”之外的其他定义（常被称为“名义定义”），希腊数学家和哲学家有着完全清楚的观念。正是在这种背景下，“存在”问题在数学中大概第一次明确地出现了。亚里士多德没有忽略指出：一个定义并不蕴涵被定义对象的存在，在此之外还必须有一个公设或一个证明。他的这个看法无疑来自数学家的实践；无论如何，欧几里得都注意在把这些概念引入论证时，公设圆的存在，并证明等边三角形、平行线、正方形等的存在（{[}153 e{]}，第 I 卷）；这些证明都是“作图”；换句话说，他以公理为依据，展示出一些数学对象，并证明它们满足他需要加以辩护的那些定义。

于是我们看到，古典时代的希腊数学达到的是一种经验性的确定性（不管它出自某位哲学家的什么样的形而上学根据）；既然不可设想推理规则会受到质疑，希腊科学的成功、再加上人们关于批判性修订并不合时宜的感觉，就在公理本身所激发的信心中占了很大一部分；这种信心近乎上个世纪人们对理论物理学诸原理所抱的那种（同样是几无限制的）信心。学园的格言“nihil est in intellectu quod non prius fuerit in sensu”（理智中无一物不是先在感觉中的）所提示的，无论如何正是这一点；而笛卡尔反对的恰恰是这句格言，认为它没有为笛卡尔希望从理性的运用中提取的东西提供足够牢固的基础。

必须等到 19 世纪初，才能看到数学家们从笛卡尔式的傲慢（更不必说康德或黑格尔的傲慢；后者对他那个时代的科学发议论未免姗姗来迟，一如时尚\(^{25}\)）中走出来，恢复到一种像希腊人那样恰当平衡的立场。对经典概念的第一次打击，是高斯、罗巴切夫斯基和鲍耶在本世纪初建立的非欧双曲几何。我们不打算在这里详细追述这一发现的来龙去脉，它是证明平行公设的无数次徒劳尝试的结果（见 {[}105 a 和 b{]}）。这一发现在当时对数学原理的影响，也许不像有时所说的那么深刻。它只是迫使人们放弃上个世纪对欧氏几何“绝对真理”的奢望，尤其是迫使人们放弃定义蕴涵公理的莱布尼茨式观点；公理不再显得全然“自明”，而是显得像一些假设，有待确定它们是否适合于对感官世界的数学表示。高斯和罗巴切夫斯基认为，各种可能几何之间的争论可以由经验来裁决（{[}206{]}，第~76~页）。黎曼的观点也是如此，他那篇著名的就职演讲《论构成几何基础的假设》旨在为各种自然现象提供一个一般的数学框架：“尚待解决的”，他说，“是这些假设在多大程度上、在什么范围内被发现为经验所证实的问题”（{[}259 a{]}，第~284~页）。但这是一个显然与数学不再有任何关系的问题；而且上述作者中似乎没有人质疑：即使一种“几何”不对应于任何实验的实在，它的定理也仍然是“数学真理”。\(^{26}\)

尽管如此，即便如此，这种信服肯定不应再归之于对经典“几何直觉”的无限制信任；黎曼试图对“n 重延展的流形”（他工作的对象）所作的描述，只是在为引入“局部坐标”寻求辩护时才依赖“直观的”\(^{27}\) 考虑；从那时起，他看来自觉已经站在了坚实的地面上，即分析的地面。但分析最终毕竟建立在实数概念之上，而这个概念直到那时还保持着非常直观的性质；函数论的进展在这个方面导致了一些非常令人不安的结果：随着黎曼本人关于积分的研究，更随着波尔查诺和魏尔斯特拉斯构造的无切线曲线的例子，数学的整个病理学开始了。一个世纪以来，我们见到的这一物种的怪物如此之多，以致有些见怪不怪了，必须把最离奇的畸形特征堆积起来，才能再让我们吃惊。但这些现象在 19 世纪大多数数学家身上引起的反应，从厌恶到惊愕不一：“直觉怎么会把我们欺骗到这种地步？”H. 庞加莱自问道（{[}251 d{]}，第~19~页）；埃尔米特则（不无一丝幽默，而这句名言的评注者们并非全都察觉到了这幽默）宣称，他“怀着恐惧与厌恶，从这场没有导数的连续函数的令人悲叹的瘟疫面前转过身去”（{[}160{]}，第 II 卷，第~318~页）。最糟糕的是，这些如此违背常识的现象，再也不能归咎于阐释得不好的观念，像“不可分量”时代那样（见第~173~页），因为它们在波尔查诺、阿贝尔和柯西的改革之后依然存在，而那次改革以与比例理论同样严格的方式确立了极限概念的基础（见第~154~页）。因此，这必须完全记在我们几何直觉的粗糙与不完备的账上；而从那以后，它作为一种证明手段被理所当然地贬黜，也是合理的。

这一认识必定会反过来作用于经典数学，而首当其冲的是几何学。无论欧几里得的公理构造曾经享有何种敬意，人们早已注意到其中不止一处缺陷，而且在古代即已如此。受到批评和证明尝试最多的，是平行公设；但欧几里得的追随者和注释者也曾试图证明其他公设（尤其是直角相等的公设），或者承认某些定义——例如直线或平面的定义——的不充分。16 世纪时，《Elements》的一位刊行者克拉维乌斯注意到缺少一条保证第四比例项存在的公设；莱布尼茨则指出，欧几里得使用几何直觉而不明确提及，例如当他承认（《Elements》第 I 卷命题 1）两个各自通过对方圆心的圆有一个公共点时（{[}198 b{]}，第 VII 卷，第~166~页）。高斯（他自己并不拒绝使用这类拓扑考虑）提请注意一个点（或一条直线）位于另外两个“之间”这一观念在欧氏作图中所起的作用，而这个观念却没有得到定义（{[}124 a{]}，第 VIII 卷，第~222~页）。最后，位移的使用——尤其是在“三角形相等的情形”中——长期被认为自明，\(^{28}\) 不久便在 19 世纪的批评者看来同样依赖于未曾陈述的公理。于是，在 1860 年至 1885 年间，出现了对几何学开端的各种局部修正（亥姆霍兹、梅雷、乌埃尔），旨在弥补其中一些缺陷。但只是到了 M. 帕施 {[}245{]}，放弃一切对直觉的诉求才成为一个正式提出并得到完全严格贯彻的纲领。他事业的成功很快为他引来了众多效法者，他们主要在 1890 年至 1910 年间，给出了彼此颇为不同的欧氏几何公理表述。这些著作中最著名的是皮亚诺用其符号语言写成的那些 {[}246 d{]}，尤其是希尔伯特的《Grundlagen der Geometrie》{[}163 c{]}，该书出版于 1899 年，以其论述的明晰与深刻，立即理所当然地成为现代公理学的宪章，甚至到了使其先行者湮没无闻的地步。的确，希尔伯特在那里不仅给出了欧氏几何的一个完备的公理系统，还把这些公理划分为不同类型的不同组，并致力于确定每一组公理的确切影响范围：不仅孤立地展开每一组的逻辑推论，而且讨论省略或修改其中某些公理所得到的各种“几何”（在这些几何中，罗巴切夫斯基几何和黎曼几何只作为特例出现\(^{29}\)）；他就这样在一个直到当时都被认为最接近感官实在的领域里，清楚地展示了数学家在选择公设时所享有的自由。尽管这些性质奇特的“元几何”使不止一位哲学家陷入窘迫，《Grundlagen》的论点还是很快被数学家们几乎一致地接受；H. 庞加莱尽管很难被指摘偏爱形式主义，也在 1902 年承认，几何公理是一些约定，对它们而言，通常所理解的“真理”概念不再有意义（{[}251 c{]}，第~66-67~页）。“数学真理”于是仅仅存在于从公理任意设定的前提出发的逻辑演绎之中。如下文将见（第~35~页以下），进行这些演绎所用的推理规则，其本身的有效性不久也将受到质疑，从而导致数学基本观念的一次彻底重塑。

